\documentclass[11pt,a4paper]{article}
\usepackage[T1]{fontenc}\usepackage[utf8]{inputenc}\usepackage{lmodern}
\usepackage[a4paper,margin=2.2cm]{geometry}
\usepackage{amsmath,amssymb}\usepackage{enumitem}\usepackage{booktabs,longtable,array}
\usepackage[hidelinks]{hyperref}
\setlength{\parindent}{0pt}\setlength{\parskip}{0.4em}
\title{\textbf{OP3 / B3: a charge-selective routing line for the
thermal-completion ledger}\\\large Recon v1 --- proposal-only; D-candidate
drafted, not inserted}
\author{Per reviewer directive after acceptance of the B2-Q1 audit}
\date{12 June 2026}
\begin{document}\maketitle

\section*{0. Scope and firewall}
B3 drafts Bridge~3 of the BB-bridge recon as \emph{one} ledger line in the
exact recorded format of the thermal-completion bookkeeping ledger
(items L1--L5: named constraint, honest caveats, ``no mechanism implied'',
strictly post-onset). The line is presented verbatim as a D-candidate; no
preprint edit is performed. Like L1--L5, it binds only after an ordered
Lorentzian branch exists; nothing is asserted at the pre-emergent level.

\section{Recorded format being matched}
The recorded ledger items each consist of: an italicised label
\emph{(Ln)~name}; a constraint statement organising any completion
mechanism; explicit caveats separating bookkeeping from prediction; and the
standing discipline that ``the mechanism \dots is not derived and is not
implied.'' The proposed line follows this pattern exactly and is intended
for insertion after \emph{(L5)} and before the Kibble--Zurek scaffold
paragraph.

\section{Proposed line (verbatim D-candidate)}
\begin{quote}
\emph{(L6) Charge-selective routing (event-ordering condition).}
Within the recorded charge assignments, the relaxing amplitude mode carries
no recorded holonomy or charge label, and within the recorded induced-gauge
architecture it supports no tree-level $\sigma F^2$ portal, loop-induced
portals reducing to the charged-sector selection rule
(\S OP-CONST block). Conversion of condensate relaxation energy into a
gauge/radiation bath therefore proceeds through the charged matter chain,
whose masses anchor to $m_f=y_fv_2$ with
$v_2=\phi_0e^{-\mathcal I_{\rm EW}}$. Any thermal-completion mechanism
hence inherits the event ordering: gradient ordering $\to$
metric/$\Phi$-sector matching (so that the $v_2$ chain exists) $\to$
charged-spectrum population $\to$ entropy production and radiation
domination (L2). This is a routing/ordering constraint on mechanisms,
conditional on the recorded charge-based selectivity; no conversion
efficiency, timescale, or mechanism is implied, and, as for L1--L5, no
statement is made at the pre-emergent level.
\end{quote}

\section{Justification and cross-checks}
\textbf{(a) Basis is preprint-recorded, not merely proposal-level.} The
charge architecture the line relies on is already in the canonical preprint
through D3/D4: the OP-CONST block records that $\sigma=\delta u$ carries no
recorded holonomy or charge label, and that the portal condition is
partially supported by the recorded induced-gauge architecture. The line
therefore cites recorded text, with the residual conditionality stated.
\textbf{(b) Same constraint class as L1--L5.} L1 bounds a temperature
ceiling, L2 a floor, L3 a duration corridor, L4 a settling tolerance, L5 an
entropy identity. The proposed L6 bounds the \emph{topology of the energy
flow}: which channel the conversion must use and in which order the
structures must exist. It adds no number and removes no mechanism freedom
beyond the routing itself.
\textbf{(c) Consistency with L4.} The $G_{\rm eff}$-settling condition and
the charged-spectrum requirement both tie to the $\Phi$-sector being
matched well before MeV scales; L6's ordering makes that shared
prerequisite explicit rather than implicit.
\textbf{(d) What it does not claim.} No efficiency, no timescale, no
mechanism; no pre-emergent statement; no claim that the ordering is
sufficient for thermal completion --- it is necessary routing only.

\section{Verdict}
\begin{enumerate}[leftmargin=1.9em,itemsep=2pt,label=\textbf{V\arabic*.}]
\item Bridge~3 is now a single, format-exact ledger-line D-candidate with a
named insertion point (after L5, before the Kibble--Zurek scaffold);
nothing has been inserted.
\item The line's basis is the recorded D3/D4 charge architecture; its
conditionality is carried inside the line text itself.
\item If accepted, the insertion is a one-line D-step with the usual
discipline (backup, gates, manifest); if declined, the BB-bridge round can
close with B1, B2 (rev 1.1), B2-Q1, and B3 all at proposal level.
\end{enumerate}
\end{document}
